\documentclass[11pt,a4paper]{article}
% Shared preamble for the three companion documents.  Deliberately minimal:
% only packages a bare TeX Live carries, since every figure is built from
% LaTeX `picture` primitives rather than from TikZ or an external plotter.
\usepackage[utf8]{inputenc}
\usepackage[T1]{fontenc}
\usepackage{amsmath,amssymb,amsthm}
\usepackage{booktabs}
\usepackage{algorithm}
\usepackage{algpseudocode}
\usepackage{xcolor}
\usepackage[margin=2.4cm]{geometry}
\usepackage[hidelinks]{hyperref}

\newcommand{\Z}{Z}
\newcommand{\Dset}{D}
\newcommand{\E}{E(P_{\Dset})}
\newcommand{\R}{\mathbb{R}}
\newcommand{\Zint}{\mathbb{Z}}
\newcommand{\fig}[1]{\input{figs/#1}}

% the legend shared by every criterion-space figure
\newcommand{\figlegend}{%
  \par\smallskip\noindent{\footnotesize
  \textcolor{green!45!black}{$\blacksquare$}~efficient point\quad
  \textcolor{black!45}{$\square$}~dominated point\quad
  \textcolor{blue!60!black}{$\square$}~the box being settled\quad
  \textcolor{red!35}{$\blacksquare$}~what the cut removes\quad
  \textcolor{red!65!black}{$\square$}~the children left\quad
  \textcolor{red!65!black}{$\circ$}~$x_B$\quad
  \textcolor{orange!85!black}{$\circ$}~the centre cut on\par}}

\newtheorem{proposition}{Proposition}
\newtheorem{remark}{Remark}


\title{Two worked examples, drawn step by step\\
\large every box, every centre, every drop --- read out of a real run}
\author{}
\date{}

\begin{document}
\maketitle

\begin{abstract}
Two instances are followed to the end, one with two variables and one with
three.  Every figure is generated by \texttt{docs/figs/regenerate.py}, which
re-runs the search and reads the boxes, maximisers and centres out of it: no
picture here was drawn by hand, and none can drift away from the code.  The
smaller example is the one the published method is stated on, so the two can be
compared line for line; the larger one is chosen because its four steps exhibit
\emph{every} branch of the algorithm --- a repair, a maximiser that is already
efficient, a box dropped on its bound, and the early exit.
\end{abstract}

\tableofcontents

\section{How to read the figures}

Each step gets one picture of \textbf{criterion space} $(\Z_1, \Z_2)$ --- not of
decision space, because that is where the algorithm keeps its unexplored
region.  In every picture:

\figlegend

\noindent
The solid blue rectangle is the box currently being settled; the dashed red
rectangles are the children it is about to be split into, drawn on the same
axes so the split can be seen rather than described.  A box open on a side
extends to the frame.

\section{Example A: the published instance, $n = 2$}

\begin{equation*}
\max \Phi(x) = \frac{x_1 + x_2 - 1}{5x_1 + x_2 - 1}
\quad\text{over}\quad
\E \ \text{of} \
\begin{cases}
  \Z_1(x) = x_1 - 2x_2 \\ \Z_2(x) = -x_1 + 4x_2
\end{cases}
\ \
\Dset :
\begin{cases}
  -2x_1 + x_2 \leqslant 0 \\
  6x_1 + x_2 \leqslant 21 \\
  -2x_1 + 4x_2 \leqslant 6 \\
  x \in \Zint^2_+
\end{cases}
\end{equation*}

$\Dset$ holds $11$ points, of which $7$ are efficient.

\begin{center}
\fig{a-decision}
\end{center}
\noindent
\textbf{Decision space.}  The eleven feasible points; the seven efficient ones
filled, the optimum of $(P_E)$ ringed.  Note that the efficient set is not a
face and not convex --- there is nothing here to hand a solver, which is the
difficulty $(P_E)$ starts from.

\paragraph{Where the answer is, and why it is not obvious.}
$\Phi$ is maximised over all of $\Dset$ at $x = (0,0)$, with $\Phi = 1$.  That
point is \emph{dominated}, so it is not a feasible answer to $(P_E)$.  The true
optimum is $\Phi = 5/17 \approx 0.294$ at $(3,3)$ --- barely a third of the
unconstrained value.  An implementation that forgot the efficiency test would
return $1$ and look entirely plausible.

\subsection*{Step 1 --- the root box}
\begin{center}\fig{a-step1}\end{center}
The box is all of $\R^2$, so the sub-problem is $\max \Phi$ over $\Dset$ itself:
$x_B = (0,0)$ with $\Phi = 1$.  The test returns $\theta > 0$ --- $(0,0)$ is
dominated --- so its value is discarded rather than banked.  The centre is the
efficient point $(2,1)$, at $\Z = (0,2)$.  $Q(2,1)$ gives $\Phi = 1/5$, the
first incumbent.  The split removes $\{\Z \leqslant (0,2)\}$ and leaves two
children: $\Z_1 \geqslant 1$, and $\Z_2 \geqslant 3$ with $\Z_1 \leqslant 0$.

\subsection*{Step 2 --- a box dropped on its bound}
\begin{center}\fig{a-step2}\end{center}
$\Z_1 \geqslant 1$.  The sub-problem cannot beat the incumbent $1/5$, so the box
is discarded \emph{whole} --- it is never split, and every descendant it would
have had is gone with it.  This is the mechanism that a good incumbent feeds,
and the reason the search is worth seeding.

\subsection*{Step 3 --- the optimum arrives, from a dominated maximiser}
\begin{center}\fig{a-step3}\end{center}
$\Z_1 \leqslant 0$, $\Z_2 \geqslant 3$.  Here $x_B = (1,2)$ with $\Phi = 1/3$ ---
larger than the incumbent, but the test rejects it: $(1,2)$ is dominated.  The
centre is $(3,3)$, and $Q(3,3) = 5/17$, which becomes the incumbent and is, in
the end, the answer.

Note what would have gone wrong twice over.  Taking $\Phi(x_B) = 1/3$ as the
incumbent would have been taking the value of a dominated point, and $1/3 >
5/17$: the cutoff would then have discarded the box holding the true optimum.

\subsection*{Steps 4 and 5 --- the rest is proof}
\begin{center}\fig{a-step4}\quad\fig{a-step5}\end{center}
Neither child can beat $5/17$; both are dropped.  The list empties and the
answer is proved optimal.  \textbf{Five sub-problems in total}, against a
feasible set of eleven points and an efficient set of seven --- the method never
enumerated $\E$, which is the whole point of it.

\section{Example B: three variables, and every branch exercised}

\begin{equation*}
\max \Phi(x) = \frac{-2x_1 + 4x_2 - 3x_3 + 1}{x_1 + 2x_2 + x_3 + 1}
\quad\text{over}\quad
\E \ \text{of} \
\begin{cases}
  \Z_1(x) = 4x_1 - x_2 + x_3 \\ \Z_2(x) = 2x_2 + 2x_3
\end{cases}
\end{equation*}
\begin{equation*}
\Dset : \quad x_1 + x_2 + 2x_3 \leqslant 5, \qquad
2x_1 + 3x_2 + x_3 \leqslant 6, \qquad
0 \leqslant x_j \leqslant 3, \quad x \in \Zint^3 .
\end{equation*}

$\Dset$ holds $15$ points; exactly $4$ are efficient, at
$\Z = (1,6), (6,4), (9,2), (12,0)$.

\begin{center}
\fig{b-decision}
\end{center}
\noindent
\textbf{Decision space, as three slices.}  With $n = 3$ the lattice is drawn
one $x_3$ at a time.  The efficient points are filled and the optimum ringed;
they lie in different slices, which is what makes the set awkward in decision
space and tidy in criterion space --- in the figures below all four sit on a
single staircase.

\subsection*{Step 1 --- repair, and the first incumbent}
\begin{center}\fig{b-step1}\end{center}
Root box.  $x_B = (0,2,0)$ with $\Phi = 9/5$, the largest value $\Phi$ takes
anywhere on $\Dset$ --- and dominated, so again not an answer.  The centre is
$(1,0,2)$ at $\Z = (6,4)$, and $Q$ there gives $\Phi = -7/4$.  That is the
incumbent: \emph{negative}, and far below the $9/5$ just seen.  The split leaves
$\Z_1 \geqslant 7$ and $(\Z_2 \geqslant 5, \Z_1 \leqslant 6)$.

\subsection*{Step 2 --- a second repair}
\begin{center}\fig{b-step2}\end{center}
$\Z_1 \geqslant 7$.  $x_B = (2,0,0)$ with $\Phi = -1$, again dominated; the
centre is $(3,0,0)$ at $\Z = (12,0)$ and $Q$ gives $-5/4$, which improves the
incumbent from $-7/4$.  Two children follow.

\subsection*{Step 3 --- a maximiser that is already efficient}
\begin{center}\fig{b-step3}\end{center}
$\Z_2 \geqslant 5$, $\Z_1 \leqslant 6$.  Here $x_B = (0,1,2)$ with
$\Phi = -1/5$ and the test returns $\theta = 0$: the maximiser \emph{is}
efficient.  So it becomes the incumbent directly, and \textbf{no $Q$ is owed} ---
$x_B$ already maximises $\Phi$ over the whole box, hence over its own slice.
This is the answer, though the search does not know that yet.

\subsection*{Step 4 --- the early exit}
\begin{center}\fig{b-step4}\end{center}
The best bound still open is $-1/5$, which the incumbent already equals, so it
cannot be beaten.  Because the open list is ordered by bound descending, no
other box can beat it either: the search stops and reports
\textbf{four boxes still open and never solved}.  Without this line all four
would have been solved, one at a time, only to be dropped.

\paragraph{What the four steps showed.}  A repair (twice), a maximiser already
efficient, a box dropped on its bound, and the early exit --- every branch of
Algorithm~1 in the companion \emph{step by step} document, on fifteen feasible
points.

\section{What the two examples say together}

\begin{table}[h]
\centering
\caption{The two runs.}
\begin{tabular}{lrrrrr}
\toprule
 & $n$ & $|\Dset|$ & $|\E|$ & sub-problems & $\Phi^{*}$ \\
\midrule
A (published instance) & 2 & 11 & 7 & 5 & $5/17$ \\
B & 3 & 15 & 4 & 4 & $-1/5$ \\
\bottomrule
\end{tabular}
\end{table}

\begin{enumerate}
\item \textbf{The unconstrained maximiser of $\Phi$ is dominated in both.}
  $\Phi = 1$ at $(0,0)$ in A against an answer of $5/17$; $\Phi = 9/5$ at
  $(0,2,0)$ in B against an answer of $-1/5$.  The efficiency test is not a
  formality on these instances --- it is the difference between an answer and a
  plausible wrong number.
\item \textbf{The optimum is reached through $Q$, not through a maximiser, in
  A.}  $(3,3)$ is found as the centre repaired to from the dominated $(1,2)$,
  and its value comes from the $Q$ sub-problem.  An implementation that skipped
  $Q$ --- reasonable-looking, if one forgets that $\Phi$ is a function of $x$ and
  not of $\Z(x)$ --- would miss it.
\item \textbf{Neither run enumerated the efficient set.}  Five sub-problems for
  seven efficient points, four for four.  On the instances of the main note the
  gap is the whole point: eight boxes settle an instance with 71 efficient
  points among 629 feasible ones.
\end{enumerate}

\end{document}
